% =====================================================================
%  Chapter 17 — The Landscape: the forced-blowup construction
% =====================================================================
\chapter{The Landscape: A Forced Blow-Up Construction, and Why It Cannot Be
De-Forced}
\label{ch:openai}

\begin{record}
\textbf{Sources of record.} \Sn{40},
\texttt{SESSION-LOG/2026-09-09-S40-openai-comparison-aniso-collapse.md}, and
\texttt{subsec:wall9-forcing} of the proof document (Wall~9, lines
10212--10254). The primary document under discussion is a 166-page writeup
accompanied by a Lean~4 formalisation repository; the session read Sections
1--3 in full, together with the repository's own summary of what the result
establishes.

\smallskip
Wall~9 is the only entry in the proof document's walls section that is
explicitly recorded \emph{outside} the mathematics of that document. It
carries no label, no part of the document uses it, and nothing in this
programme rests on it. It is there to bound what future sessions may import.
\end{record}

In September 2026, while this programme was working on the annulus residual, a
finite-time blow-up construction for the \emph{forced} three-dimensional
incompressible Navier--Stokes equations was published. The owner asked two
questions: does it bear on what T-First is actually trying to do, and if there
is anything usable in it, build that --- explicitly without introducing a
forcing term, since a forcing term leaves the Prize equations.

This chapter is the answer to both. The short form: it is a real result on the
problem it addresses, it addresses a different problem, and its mechanism
cannot be transplanted. What could be taken was taken, and is measured in
\S\ref{sec:aniso}.

\section{The four alternatives}
\label{sec:four-alternatives}

Fefferman's official statement of the Millennium problem offers four
alternatives, and a proof of any one of them resolves it. Two are positive
statements about smoothness; two are negative statements about breakdown. They
differ along a second axis as well, and that second axis is the whole subject
of this chapter.

\begin{figure}[ht]
\centering
\begin{tikzpicture}[
  cell/.style={draw, rounded corners=2pt, align=center, inner sep=5pt,
               font=\small, text width=44mm, minimum height=17mm},
  hit/.style={cell, very thick, draw=impossiblered, fill=impossiblered!7},
  ours/.style={cell, very thick, draw=provedgreen, fill=provedgreen!7},
  hdr/.style={font=\small\bfseries, align=center}
]
\node[hdr] at (-3.6, 2.3) {};
\node[hdr] at (0.0, 2.3) {on \(\R^3\)};
\node[hdr] at (5.2, 2.3) {on \(\T^3=\R^3/\Z^3\)};

\node[hdr, align=right, text width=30mm] at (-3.9, 0.9)
  {\textbf{(A), (B)}\\ existence and\\ smoothness};
\node[hdr, align=right, text width=30mm] at (-3.9,-1.6)
  {\textbf{(C), (D)}\\ breakdown};

\node[ours] (A) at (0.0, 0.9)
  {\textbf{(A)}\\ \(f\equiv0\)\\[2pt]
   \footnotesize\emph{the T-First target}};
\node[ours] (B) at (5.2, 0.9)
  {\textbf{(B)}\\ \(f\equiv0\)\\[2pt]
   \footnotesize\emph{the T-First target}};
\node[hit] (C) at (0.0,-1.6)
  {\textbf{(C)}\\ forcing \(f\) permitted\\[2pt]
   \footnotesize\emph{established}};
\node[hit] (D) at (5.2,-1.6)
  {\textbf{(D)}\\ forcing \(f\) permitted\\[2pt]
   \footnotesize\emph{established, via compact support}};
\end{tikzpicture}
\caption[The four Clay alternatives]{Fefferman's four alternatives. The
positive statements (A) and (B) are posed with \(f\equiv0\); the breakdown
statements (C) and (D) admit an external forcing term subject to smoothness
and decay conditions. The construction discussed here establishes (C) and, via
its compact-support cutoff, (D). This programme works on (A) and (B), and has
proved neither.}
\label{fig:four-alternatives}
\end{figure}

The writeup states outright which alternatives it establishes, and the
programme's record repeats it without gloss: \emph{in Fefferman's formulation
the result establishes alternatives (C) and (D), not (A)/(B)}. That is the
proof document's sentence, at \texttt{subsec:wall9-forcing}, and it is the one
this book uses.\footnote{The \Sn{40} session log goes one step further and
describes (A)/(B) as ``the unforced conditions the \$1M Prize covers''. The
proof document does not say that, and by the source hierarchy of
\S\ref{sec:source-hierarchy} the document governs: it records which
alternatives are established and stops there. The stronger claim about prize
eligibility is a reading of the official problem statement, not a finding of
this programme, and this book does not adopt it.}

\section{What the construction establishes}
\label{sec:their-theorem}

The main theorem, as recorded in \Sn{40}:

\begin{record}
\textbf{Theorem 1.1 (as stated in the source).} For every \(\nu>0\) there
exist a smooth, compactly supported (in space and time) forcing \(f\), a
compact set \(K\), and smooth \((u,p)\) on \(\R^3\times[0,1)\) solving the
\emph{forced} Navier--Stokes equations with \(u(\cdot,0)=0\), such that the
kinetic energy is uniformly bounded while
\[
  \limsup_{t\uparrow1}\ \norm{u(t)}_{L^\infty} \;=\; \infty .
\]
\end{record}

Two features of the statement should be noted before the mechanism, because
they are what make it a strong result rather than a soft one. The viscosity is
\emph{arbitrary and fixed}: the construction is not a small-viscosity or
inviscid-limit argument. And the initial datum is \emph{zero}: the entire
singular structure is generated from rest by the forcing within a finite time,
inside a compact set.

\section{How it works}
\label{sec:their-mechanism}

The mechanism is an axisymmetric, \emph{anisotropic}, self-similar collapsing
vortex core. Write \(\tau=1-t\) for the time to the singularity. The core has
two length scales, and they collapse at different rates:
\[
  \ell_r \asymp \tau^{1/2},
  \qquad
  \ell_z \asymp \tau^{1/2-h},
  \qquad 0<h<\tfrac1{100}.
\]
The radial scale collapses at exactly the diffusive rate; the axial scale
collapses \emph{more slowly}, by the small exponent \(h\). The consequence is
the design feature of the whole construction: the two Reynolds numbers
separate,
\[
  \mathrm{Re}_\theta \asymp \tau^{-h}\to\infty,
  \qquad
  \mathrm{Re}_r = O(1).
\]
Viscosity remains in full control radially --- which is why the collapse is not
simply diffused away --- while only the azimuthal, spin-up direction escapes
it. That is a genuinely clever way around the fact that at the diffusive
scaling, viscosity ordinarily wins.

\begin{figure}[ht]
\centering
\begin{tikzpicture}[>={Latex[length=2mm]}]
  % axis
  \draw[thick, gray] (0,-2.5) -- (0,2.5) node[above, font=\scriptsize]{\(z\)};
  \draw[->, thick, gray] (-3.4,0) -- (3.4,0) node[right, font=\scriptsize]{\(r\)};

  % early core (wide)
  \draw[openblue, thick, dashed] (0,0) ellipse (2.3 and 1.9);
  \node[font=\scriptsize, openblue] at (2.55,1.55) {\(t\) early};

  % late core (tall thin)
  \draw[impossiblered, very thick] (0,0) ellipse (0.55 and 1.55);
  \node[font=\scriptsize, impossiblered] at (1.05,-1.75) {\(t\to1\)};

  % scale annotations
  \draw[<->, thick] (-0.55,-1.95) -- node[below, font=\scriptsize]
    {\(\ell_r\asymp\tau^{1/2}\)} (0.55,-1.95);
  \draw[<->, thick] (1.55,-1.55) -- node[right, font=\scriptsize]
    {\(\ell_z\asymp\tau^{1/2-h}\)} (1.55,1.55);

  % swirl arrows
  \foreach \y in {-0.8,0,0.8} {
    \draw[->, provedgreen, thick] (-0.9,\y) arc (180:0:0.9 and 0.22);
  }
  \node[font=\scriptsize, provedgreen] at (-2.15,0.35)
    {swirl: \(\mathrm{Re}_\theta\!\to\!\infty\)};
  \node[font=\scriptsize, gray, align=center] at (-2.35,-1.5)
    {radial:\\ \(\mathrm{Re}_r=O(1)\)\\ viscosity in control};

  % pulses
  \foreach \y in {1.15,-1.15} {
    \draw[decorate, decoration={snake, amplitude=0.7mm, segment length=2.4mm},
          withdrawnorange, thick] (-1.9,\y) -- (1.9,\y);
  }
  \node[font=\scriptsize, withdrawnorange, align=center] at (2.75,1.15)
    {ring pulses};
\end{tikzpicture}
\caption[The anisotropic collapsing core]{The geometry: an axisymmetric vortex
core whose radial and axial scales collapse at different rates. Because
\(\ell_z\) collapses more slowly than \(\ell_r\), the azimuthal Reynolds number
diverges while the radial one stays \(O(1)\), so viscosity controls the radial
direction and only the swirl escapes. The oscillatory ring pulses supply the
mean momentum flux that cancels the leading singular part of the residual.}
\label{fig:aniso-geometry}
\end{figure}

The base collapsing profile is not by itself a solution of anything: as
\(t\to1\) it leaves an unbounded momentum residual. The construction's
substantive work is in cancelling that residual. Spatially oscillatory,
ring-shaped ``pulses'' are introduced whose \emph{nonlinear self-advection}
generates a mean momentum flux that cancels the leading singular part;
successive corrections handle the remainder; and a compact-support cutoff
localises the whole object, which is what upgrades (C) to (D).

This is a direct extension of the vortex-layer-amplification programme of
C\'ordoba and Mart\'inez-Zoroa, cited as such in the source. Stated plainly,
without disparagement: the construction solves a hard problem, and the
anisotropic scale separation and the pulse-cancellation device are real
contributions.

\section{The de-forcing question}
\label{sec:de-forcing}

The programme's question was not whether the result is correct. It was:
\emph{can the construction be run with \(f=0\)?} If it could, it would be a
statement about the Prize equations.

It cannot, and the reason is worth stating precisely, because ``it uses a
forcing term'' is not by itself an argument --- many constructions use a
forcing term that turns out to be removable or absorbable.

\begin{record}
\textbf{The forcing is load-bearing at three separate points, not
decorative.}
\begin{enumerate}[label=(\alph*), leftmargin=2.2em]
\item An exponentially small external force \emph{seeds each oscillatory
  pulse}. The pulses do not arise from the dynamics; they are injected.
\item After pulse cancellation, \emph{what remains of the momentum residual
  becomes part of \(f\)}.
\item The compact-support cutoff introduces \emph{its own} force term.
\end{enumerate}

\smallskip
And the structural fact that settles it: \(f\) is \emph{defined} as the
Navier--Stokes residual \(R(u,p)\) of the constructed field.
\end{record}

That last sentence is the whole argument. The construction proceeds by writing
down a velocity field with the desired singular behaviour and then declaring
the forcing to be whatever the equations fail to satisfy. There is therefore
nothing to remove: setting \(f\equiv0\) is not a simplification of the
construction, it is the assertion that this exact profile solves the unforced
equations --- which is a different question, which the source does not claim to
answer, and which the source states it does not establish.

\begin{record}
\textbf{Wall 9 (\texttt{subsec:wall9-forcing}).}
\emph{Classification:} \IMPOSSIBLE{} \emph{as an import route} --- there is
nothing to de-force, since the residual is the forcing by construction. The
underlying blow-up question for the unforced equations is of course \OPEN{};
that is the Prize.

\smallskip
\emph{Forbidden:} forcing terms, modified equations and ``weak internal
forces'' as mechanisms anywhere in this programme; reading a forced blow-up
result as evidence about the unforced equations \emph{in either direction}.

\smallskip
\emph{Permitted:} the \emph{geometry} of such constructions as a source of
adversarial initial data for the unforced solver.
\end{record}

The prohibition runs in both directions, and this matters. It would be equally
illegitimate to read the construction's existence as evidence that unforced
blow-up is likely, and to read the difficulty of de-forcing it as evidence that
unforced blow-up is unlikely. Neither inference is available.

One substantive observation was recorded, and recorded \emph{as} an
observation rather than as evidence: even to make a fixed-viscosity velocity
field blow up, the construction needed an externally injected, precisely tuned
force to sustain the anisotropic collapse against dissipation. The collapse
does not sustain itself under real viscosity alone. That is a data point
consistent with the general intuition that dissipation opposes this particular
collapse mode absent external forcing --- and no more than that.

Two further factual findings from the reading, recorded for completeness. The
document contains no mention of temperature or thermodynamics anywhere: a full
grep of the 166 pages returned zero hits for \emph{temperature}, \emph{thermal}
and \emph{thermodynam}, and \(\nu\) is a fixed constant throughout. There is
no structural overlap with Route~1's \(\mu(T)\), \(k(T)\), \(A(T)\) framework
of Chapter~\ref{ch:origin}. And the construction has no bearing on any open
item in the \Sn{31}--\Sn{39} programme: not the annulus step, not
\(\sstar\)-decay, not H2, not hypothesis~(R), not gap~(A-up).

\section{What the programme took anyway}
\label{sec:aniso}

The permitted import is the geometry, and it was taken under a standing rule
of the owner's, stated at the time the work was commissioned: \emph{inspiration
yes, forcing term never, because a forcing term leaves the Prize equations.}

\subsection{The initial condition}

\texttt{layer3/aniso\_collapse\_IC.py} builds an axisymmetric,
independently-tunable-anisotropy, swirl-plus-outflow vortex column:
\[
  u_\theta(r,z) = \frac{r}{\ell_r}\,
    \exp\!\Bigl(-\frac{r^2}{\ell_r^2}-\frac{z^2}{\ell_z^2}\Bigr),
  \qquad
  u_z(r,z) = \Bigl(\frac{z}{\ell_z}+b\Bigr)
    \exp\!\Bigl(-\frac{r^2}{\ell_r^2}-\frac{z^2}{\ell_z^2}\Bigr),
\]
with \(\ell_z=\text{(aspect ratio)}\times\ell_r\) and a small asymmetric axial
bias \(b\), mirroring the source's deliberate breaking of exact
\(z\to-z\) reflection symmetry. The radial component and exact
divergence-freeness are recovered by the existing spectral Leray projection.
The amplitude is calibrated so that the total kinetic energy \emph{matches} the
existing adversarial self-similar profile at the same resolution --- so the new
test is not weaker than the old one by construction.

Four properties of the build deserve emphasis, because together they are what
make it a legitimate import rather than a smuggled forcing term:

\begin{enumerate}[leftmargin=2.2em]
\item The initial condition is fed to the \emph{existing, completely
  unmodified} \texttt{route2\_3D\_jax} solver, which carries \(f\equiv0\)
  throughout --- the same solver used by every other Layer-3 Route-2
  experiment.
\item No forcing term is introduced anywhere in the module.
\item The module does \emph{not} reproduce the source's analytic profile,
  which would require solving its appendix ODE and analytic-continuation
  system. It is a qualitative surrogate sharing the one structural feature ---
  independently tunable radial and axial anisotropy --- that the construction
  exploits.
\item \(31\) tests, all passing, covering shapes, energy calibration equality
  across aspect ratios, \(\ell_z\) scaling and its grid floor, rejection of
  invalid aspect ratios, post-Leray divergence, and short-run positivity of
  the measured exponents. No existing file was edited; the \(48\)-test
  self-similar-profile suite was re-run unmodified and still passes.
\end{enumerate}

\subsection{What the sweep returned}

\begin{table}[ht]
\centering
\small
\begin{tabular}{@{}l r r r r r l@{}}
\toprule
\texttt{exp\_id} & \(N\) & aspect & floor-clamped? & steps &
\(\delta_{\max}\) & verdict \\
\midrule
\texttt{EXP-L3-R2-ANISO1000-JAX-064} & \(64\) & \(1.000\) & no & \(20\) & \(1.50\) & \textsf{PASS} \\
\texttt{EXP-L3-R2-ANISO0500-JAX-064} & \(64\) & \(0.500\) & no & \(20\) & \(1.50\) & \textsf{PASS} \\
\texttt{EXP-L3-R2-ANISO0200-JAX-064} & \(64\) & \(0.200\) & no & \(26\) & \(0.50\) & \textsf{PASS} \\
\texttt{EXP-L3-R2-ANISO0100-JAX-064} & \(64\) & \(0.100\) & yes & \(29\) & \(0.50\) & \textsf{PASS} \\
\texttt{EXP-L3-R2-ANISO0100-JAX-128} & \(128\) & \(0.100\) & no & \(64\) & \(0.50\) & \textsf{PASS} \\
\texttt{EXP-L3-R2-ANISO0050-JAX-128} & \(128\) & \(0.050\) & \textbf{yes} & \(69\) & \(\mathbf{0.00}\) & \textsf{FAIL} \\
\bottomrule
\end{tabular}
\caption[The anisotropic-collapse sweep]{All six rows logged under
\texttt{claim\_id}~\texttt{claim\_l3\_aniso\_collapse} in
\texttt{route2\_jax\_experiments}. \(\nu=10^{-3}\), \(\varepsilon=0.1\),
\(t_{\mathrm{end}}=1.0\). Energy decayed in every run and \(\theta\) stayed
non-negative in every run, including the failure.}
\label{tab:aniso}
\end{table}

The result requires care, and \Sn{40} took it. Two things are visible.

\emph{A real trend.} \(\delta_{\max}\) weakens monotonically as anisotropy
increases: \(1.50,1.50,0.50,0.50\). That is not noise, and it was flagged as
worth tracking.

\emph{One failure, and a control that explains it.} The single \textsf{FAIL},
at aspect ratio \(0.05\) and \(N=128\), returned \(\delta_{\max}=0.00\)
exactly. It is not a truncated or vacuous run: \(\theta\) stayed non-negative,
the energy decayed normally from \(1.328\times10^{-3}\) to
\(7.63\times10^{-4}\), and it completed all \(69\) steps to
\(t_{\mathrm{final}}=1.0\). But before treating it as a finding, the session
checked for a discretisation artefact --- and found one. At \(N=128\), the
requested \(\ell_z=0.05\times\ell_r=0.0393\) falls \emph{below one grid cell},
so the module clamps it to its hard floor of \(2\pi/128=0.0491\): the core's
axial extent is representable as literally one grid point.

The control run isolates it. At aspect ratio \(0.100\), \(N=128\), the axial
scale is \(0.0785\) and is \emph{not} clamped; that run returns
\(\delta_{\max}=0.50\), identical to the \(N=64\) result at the same aspect
ratio. So aspect ratio \(0.1\) is genuinely resolution-converged, and the sole
failure coincides exactly with the one configuration pinned to the grid floor.

\Sn{40}'s own statement of the conclusion is the one this book adopts, and it
is deliberately weaker than the evidence might tempt: this is inference from a
single controlled comparison, not proof that the floor explains everything.
Testing the real aspect-ratio-\(0.05\) geometry needs \(N\) large enough for
\(\ell_z\) to clear the floor by several cells --- \(N\gtrsim512\) --- which was
beyond the session's compute budget and was not attempted. The honest state of
the sweep is: \(\delta_{\max}>0\) at every properly resolved aspect ratio
tested, down to \(0.1\), confirmed at two resolutions; \emph{no genuine
adversarial anisotropy threshold has been found, only a hard numerical
resolution floor}.\footnote{A small internal inconsistency in the source, noted
for anyone reconciling the record: the \Sn{40} log's ``Files changed'' section
reports \(4\) new rows in \texttt{route2\_jax\_experiments}, written before the
same session's follow-up added the two \(N=128\) runs. The database has six
rows carrying \texttt{claim\_l3\_aniso\_collapse}, and the database governs.}

The epistemic status is the same as that of the M9 blow-up search of
Chapter~\ref{ch:layer3}: finite \(N\) cannot resolve an arbitrarily thin core,
and a pass down to aspect ratio \(0.1\) does not rule out failure at ratios
thinner than the grid resolves.

\section{Their setting against ours}
\label{sec:comparison}

\begin{table}[ht]
\centering
\small
\begin{tabular}{@{}p{0.19\textwidth} p{0.35\textwidth} p{0.35\textwidth}@{}}
\toprule
& \textbf{The forced construction} & \textbf{T-First} \\
\midrule
Equations
  & Forced incompressible Navier--Stokes, \(\partial_tu+(u\!\cdot\!\nabla)u
    =-\nabla p+\nu\Delta u+f\)
  & The exact Prize equations, \(f\equiv0\), from \Sn{18} onward \\[3pt]
Forcing
  & Essential. Smooth, compactly supported in space and time; \emph{defined}
    as the Navier--Stokes residual of the constructed field
  & Forbidden by standing rule; every solver in the programme carries
    \(f\equiv0\) \\[3pt]
Domain
  & \(\R^3\), with a compact-support cutoff giving \(\T^3\) as well
  & \(\T^3=\R^3/(2\pi\Z)^3\) throughout \\[3pt]
Symmetry
  & Axisymmetric with swirl, with \(z\to-z\) deliberately broken
  & General, no symmetry class assumed; \Sn{46} proposed an
    axisymmetric-with-swirl pivot, not yet executed \\[3pt]
Viscosity
  & Arbitrary fixed \(\nu>0\); not a small-viscosity argument
  & Arbitrary fixed \(\nu>0\) in Route~2; \(\mu(T)\) in Route~1 \\[3pt]
Initial data
  & \(u(\cdot,0)=0\); the structure is generated from rest by \(f\)
  & \(u_0\in C^\infty(\T^3)\) or \(H^1(\T^3)\); nothing is injected after
    \(t=0\) \\[3pt]
Direction of the result
  & \textbf{Construct} a singularity
  & \textbf{Exclude} a singularity \\[3pt]
Clay alternative
  & (C), and via compact support (D)
  & (A) and (B) \\[3pt]
What success wins
  & A finite-time blow-up theorem for the forced equations; the source states
    it establishes (C)/(D) and does not claim (A)/(B)
  & Global regularity for the unforced equations --- the Prize. Not achieved:
    the claim of \Sn{29} was withdrawn in \Sn{31} \\[3pt]
Status
  & A completed construction, with a Lean~4 formalisation effort
  & \OPEN. The strongest unconditional results are Claim~A and regularity for
    shear and near-shear data \\[3pt]
Overlap with this programme
  & None analytically: no temperature, no thermodynamics, no bearing on any
    open item of \Sn{31}--\Sn{39}
  & One import: the collapse \emph{geometry}, as an unforced adversarial
    initial condition \\
\bottomrule
\end{tabular}
\caption[The two settings compared]{Row by row. The two lines of work address
different alternatives of the same problem statement, in opposite directions,
with the forcing term as the dividing line.}
\label{tab:comparison}
\end{table}

\section{A note on how it was reported}
\label{sec:reporting}

One factual caution, recorded because a reader of this book may encounter the
secondary accounts before the primary source.

The construction was widely described in public as a solution of the
Navier--Stokes Millennium Prize problem. The distinction that account loses is
the one this chapter has been about: the theorem proved is a finite-time
blow-up statement for the \emph{forced} equations, establishing Fefferman's
alternatives (C) and (D), with a forcing term that is defined as the residual
of the constructed field. The source itself states which alternatives it
establishes, and does not claim the unforced result.

\Sn{40} also records that the public discussion carried a credit dispute,
involving separately posted forced-blow-up results for the incompressible
porous-medium equation, two-dimensional Boussinesq and three-dimensional Euler
extending the same C\'ordoba--Mart\'inez-Zoroa programme, and that those
preprints are not acknowledged in the writeup the session read. That is
context; it is not mathematics, and nothing in this book turns on it.

The programme's own position is the one stated in Wall~9 and repeated here
without amplification: a forced blow-up result is not evidence about the
unforced equations in either direction, the geometry may be borrowed, and the
forcing may not.
